\exercisegroup

\begin{enumerate}
\def\labelenumi{\arabic{enumi})}
\tightlist
\item
设 S 是从 ]0, 1[ 到 R 的映射 \(t \mapsto \sin(\pi/t)\) 的图。设 W 是 S 与一条位于 \(\mathbf{R}^2 - \mathbf{S}\) 中、连接点 (0, 0) 和 (1, 0) 的道路的像的并（“Warsaw circle”）。
\end{enumerate}

\begin{enumerate}
\def\labelenumi{\alph{enumi})}
\item
证明空间 W 是单道路连通的，但不是单连通的。
\item
对于每个整数 \(n \geqslant 0\)，置 \(X_{n} = R/2^{n}Z\)，并记 \(f_{n}: X_{n+1} \to X_{n}\) 为典范满射；令 X 为 \(\varprojlim_{n} X_{n}\) 所成的空间。证明空间 X 是单道路连通的，但不是单连通的。
\end{enumerate}

\begin{enumerate}
\def\labelenumi{\arabic{enumi})}
\setcounter{enumi}{1}
\item
设 X 是 III, p. 331 的习题 4 中的拓扑空间，令 Y 为它的悬挂（I, p. 149, exerc. 1）。证明空间 Y 是单连通的，但 \(\pi_1(Y, y) \neq \{e_y\}\) 对所有 \(y \in Y\) 都成立。
\item
设 X 是一个拓扑空间，x 是 X 的一个点。
\end{enumerate}

设 R 是 X 的覆盖范畴，E 是集合范畴（cf. II, p. 160, example 2）。

\begin{enumerate}
\def\labelenumi{\alph{enumi})}
\item
假设对于 R 的每个对象 E，E 这个 X 的覆盖的纤维 \(E_{x}\) 是 E 的一个对象；于是置 \(\Phi_{x}(\mathrm{E}) = \mathrm{E}_{x}\)。对于 R 中作为对象的 B 的覆盖之间的每个态射 \(f: E' \to E\)，记 \(\Phi_{x}(f)\) 为 f 在纤维上诱导的从 \(E'_{x}\) 到 \(E_{x}\) 的映射。证明 \(\Phi_{x}\) 是从范畴 R 到范畴 E 的函子。
\item
令 \(\mathrm{G}(x)\) 为由族 \(\varphi = (\varphi_{\mathrm{E}})\) 构成的集合，其中对于 R 的每个对象 E，\(\varphi_{E}\) 是集合 \(\Phi_{x}(\mathrm{E})\) 的一个置换，并且对于 R 的每个箭头 \(f: E' \to E\)，有 \(f_{*} \circ \varphi_{E'} = \varphi_{E} \circ f_{*}\)（“函子 \(\Phi_{x}\) 的自同构”）。
\end{enumerate}

由关系 \(\varphi \cdot \psi = (\varphi_{\mathrm{E}} \circ \psi_{\mathrm{E}})\) 对于 \(\varphi, \psi \in \mathrm{G}(x)\)，赋予集合 \(\mathrm{G}(x)\) 群结构。对于 \(\mathcal{R}\) 的每个对象 E，记 \(\mathrm{G}(x)_{\mathrm{E}}\) 为所有 \(\varphi \in \mathrm{G}(x)\) 且满足 \(\varphi_{\mathrm{E}} = \mathrm{Id}_{\Phi_x(\mathrm{E})}\) 的元素所成的集合；它是 \(\mathrm{G}(x)\) 的子群。

\begin{enumerate}
\def\labelenumi{\alph{enumi})}
\setcounter{enumi}{2}
\item
证明存在 \(\mathrm{G}(x)\) 上唯一的拓扑群结构，使得各群 \(\mathrm{G}(x)_{\mathrm{E}}\) 在单位元处构成邻域基。
\item
设 T 是 X 的单连通子空间。若 x 和 y 是 T 的点，则拓扑群 \(\mathrm{G}(x)\) 和 \(\mathrm{G}(y)\) 同构。
\item
证明道路的提升定义了从 \(\pi_{1}(X,x)\) 到 \(G(x)\) 的连续群同态。这个同态不必是满射（exerc. 1; I, p. 146, exerc. 6），也不必是单射（exerc. 2）。若 X 可解圈，且 X 的每个连通覆盖都同构于 R 的一个对象，则它是一个同构。
\end{enumerate}

\begin{enumerate}
\def\labelenumi{\arabic{enumi})}
\setcounter{enumi}{3}
\tightlist
\item
对于每个整数 \(n \geqslant 1\)，令 \(X_{n}\) 为坐标平面 \(R^{2}\) 中三个线段的并，它们的端点是 \((0,1)\)、\((1/2n,0)\) 和 \((1/(2n-1),0)\)。令 \(X_{0}\) 为端点是 \((0,1)\) 和 \((0,0)\) 的线段。令 X 为集合 \(X_{n}\)（对于 \(n \geqslant 1\)）的并，并令 Y = \(\overline{X}\)。令 \(a = (0,1)\)。
\end{enumerate}

\begin{enumerate}
\def\labelenumi{\alph{enumi})}
\item
证明 \(Y = X \cup X_{0}\)。
\item
证明 X 是局部可缩的。
\item
证明 X 到 Y 的典范嵌入诱导了一个从 \(\pi_1(\mathrm{X},a)\) 到 \(\pi_1(\mathrm{Y},a)\) 的群同构。
\end{enumerate}

\(d)\) 证明 \(\pi_1(Y,a)\) 上的紧致收敛拓扑的商拓扑不是离散的。

\begin{enumerate}
\def\labelenumi{\arabic{enumi})}
\setcounter{enumi}{4}
\tightlist
\item
设 X 是 \(\mathbf{R}^2\) 的子空间，它是 \([0,1] \times \{0,1\}\)、\(\{0\} \times [0,1]\) 以及 \(\{\frac{1}{n}\} \times [0,1]\)（\(n \geqslant 1\)）的并。
\end{enumerate}

\begin{enumerate}
\def\labelenumi{\alph{enumi})}
\item
证明 X 的每一点 a 都有一个邻域 V，使得 \(\pi_{1}(V,a)\) 在 \(\pi_{1}(X,a)\) 中的像是平凡的。
\item
证明群 \(\pi_{1}(X,a)\) 上的紧致收敛拓扑的商拓扑不是离散的。
\end{enumerate}

\begin{enumerate}
\def\labelenumi{\arabic{enumi})}
\setcounter{enumi}{5}
\tightlist
\item
令 X 为 \(R^{2}-Q^{2}\) 这个空间。
\end{enumerate}

\begin{enumerate}
\def\labelenumi{\alph{enumi})}
\item
证明空间 X 是道路连通的。
\item
设 \(a, b\) 是 \(\mathbf{R} - \mathbf{Q}\) 中的点。令 \(c_{a,b}: \mathbf{I} \to \mathbf{R}^2\) 为满足 \(c(0) = (a, a)\)、\(c(1/4) = (a, b)\)、\(c(1/2) = (b, b)\)、\(c(3/4) = (b, a)\) 以及 \(c(1) = (a, a)\) 的唯一映射，并且它在区间 \([(m-1)/4, m/4]\) 上的限制对于 \(1 \leqslant m \leqslant 4\) 是仿射的。证明 \(c_{a,b}\) 是 \(\mathbf{X}\) 中的圈。
\item
设 \(a, b, b'\) 是 \(\mathbf{R} - \mathbf{Q}\) 中的元素，且 \(b \neq b'\)。证明圈 \(c_{a,b}\) 和 \(c_{a,b'}\) 不严格同伦。
\end{enumerate}

\(d)\) 推出，对于每个 \(a \in \mathbf{R} - \mathbf{Q}\)，群 \(\pi_1(\mathrm{X}, (a, a))\) 的基数为连续统。
% semantic-fragments: legacy-input-seam
\relax{}
